\documentclass[11pt,a4paper,openany]{book}

\makeatletter\def\input@path{{../_shared/}}\makeatother
\usepackage{course}

\begin{document}
\raggedbottom

\coursetitle{Modèles stochastiques pour les communications}{BA5}{2026--2027}

\tableofcontents
\clearpage

% =========================================================
\chapter{Module 1 : Variables aléatoires}
% =========================================================

\begin{rembox}[title=À propos de ces notes]
Seuls les points nécessitant une explication plus poussée sont détaillés ici. Le reste
(définitions, formules, propriétés) se trouve déjà dans la cheatsheet.
\end{rembox}

\section{Variable aléatoire}

\subsection{Rappels}

\begin{defnbox}[title=Fonction de répartition]
Une v.a. est caractérisée par sa \textbf{fonction de répartition} $F_X(x)$, qui est la
probabilité que $X$ prenne une valeur inférieure ou égale au réel $x$ :
\[
F_X(x) = P(X \leq x) = P(A), \qquad A = \{\zeta \in \Omega \mid X(\zeta) \leq x\}.
\]
\end{defnbox}

\begin{defnbox}[title=Densité de probabilité]
La fonction de \textbf{densité de probabilité} est définie par
\[
f_X(x) = \frac{dF_X(x)}{dx}.
\]
\end{defnbox}

\subsection{Fonction d'une variable aléatoire}

Soit $g(\cdot)$ une fonction continue. Si $X$ est une v.a. continue, l'expression
$Y = g(X)$ désigne également une v.a. continue. Soit $y$ une valeur particulière de $Y$, et
$x_1, x_2, \ldots, x_m$ les $m$ racines respectives de l'équation $y = g(x)$.

\begin{propbox}[title={Densité de $Y = g(X)$}]
La densité de probabilité de $Y = g(X)$ est donnée par
\[
f_Y(y) = \sum_i \frac{f_X(x_i)}{|g'(x_i)|}
\]
où la somme porte sur toutes les racines $x_i$ de $y = g(x)$. Cette formule est tout à fait
générale, et s'applique à tous les cas où $X$ est continue.
\end{propbox}

\subsection{Espérance d'une fonction d'une variable aléatoire}

\begin{defnbox}[title=Espérance]
L'\textbf{espérance} de la fonction $g(X)$, où $X$ est une v.a. de densité de probabilité
$f_X(x)$, est définie par
\[
\E[g(X)] = \int_{-\infty}^{\infty} g(x) f_X(x)\,dx
\]
pour autant que cette intégrale converge absolument.
\end{defnbox}

\begin{defnbox}[title=Moment d'ordre $n$]
Si $g(X) = X^n$, l'espérance porte le nom de \textbf{moment d'ordre $n$} :
\[
\E[X^n] = \int_{-\infty}^{\infty} x^n f_X(x)\,dx.
\]
\end{defnbox}

\begin{defnbox}[title=Variance et écart-type]
La v.a. $Y = X - \mu_X$ est une v.a. \textbf{centrée} vu que $\E[Y] = 0$. La
\textbf{variance} de la v.a. $X$ est le moment d'ordre 2 de $Y = X - \mu_X$, et l'écart-type
$\sigma_X$ est la racine carrée de la variance :
\begin{align*}
\Var[X] = \sigma_X^2 &= \E[(X - \mu_X)^2] = \int_{-\infty}^{\infty} (x - \mu_X)^2 f_X(x)\,dx
\\
&= \E[X^2] - \mu_X^2 = \int_{-\infty}^{\infty} x^2 f_X(x)\,dx - \mu_X^2.
\end{align*}
Dans le cas discret, avec $P(X = x_i) = p_i$,
\[
\mu_X = \E[X] = \sum_i p_i x_i, \qquad \Var[X] = \sigma_X^2 = \E[(X-\mu_X)^2] = \sum_i p_i
(x_i - \mu_X)^2.
\]
\end{defnbox}

\begin{exbox}[title=Une v.a. sans moments -- la loi de Cauchy]
Les moments peuvent ne pas exister pour certaines v.a. La densité de probabilité de la v.a.
de Cauchy est
\[
f_X(x) = \frac{1}{\pi (1 + x^2)}.
\]
On vérifie qu'il s'agit bien d'une densité de probabilité :
\[
\int_{-\infty}^{\infty} f_X(x)\,dx = \frac{1}{\pi}\Big[\arctan(x)\Big]_{-\infty}^{\infty} = 1.
\]
Par contre,
\[
\int_{-\infty}^{\infty} |x f_X(x)|\,dx = \frac{2}{\pi} \int_0^{\infty} \frac{x}{1+x^2}\,dx =
\frac{1}{\pi}\Big[\ln(1+x^2)\Big]_0^{\infty} = \infty,
\]
donc $x f_X(x)$ n'est pas absolument intégrable, et $\E[X]$ n'est pas défini.
\end{exbox}

\subsection{Transformées}

\begin{defnbox}[title=Fonction caractéristique]
L'opérateur d'espérance mathématique appliqué à la v.a. complexe $\exp(j\omega X)$, où $j$
est la racine carrée de $-1$, définit la \textbf{fonction caractéristique} $\Phi_X(\omega)$
de la v.a. $X$ :
\[
\Phi_X(\omega) = \E[e^{j\omega X}] = \int_{-\infty}^{\infty} e^{j\omega x} f_X(x)\,dx.
\]
C'est le complexe conjugué de la transformée de Fourier de $f_X(x)$. Elle existe toujours car
$|\Phi_X(\omega)| \leq \int_{-\infty}^{\infty} |f_X(x)|\,dx = 1$, et est maximale en $\omega =
0$ car $\Phi_X(0) = 1$.
\end{defnbox}

\begin{propbox}[title=Transformée inverse]
\[
f_X(x) = \frac{1}{2\pi} \int_{-\infty}^{\infty} e^{-j\omega x} \Phi_X(\omega)\,d\omega.
\]
Si deux v.a. continues ont la même fonction caractéristique, elles ont la même densité de
probabilité, et vice-versa.
\end{propbox}

\begin{propbox}[title=Moments via la fonction caractéristique]
En dérivant $\Phi_X(\omega)$ par rapport à $\omega$ et en évaluant en $\omega = 0$ :
\[
\frac{d\Phi_X(\omega)}{d\omega}\bigg|_{\omega=0} = j \E[X] \quad \Longrightarrow \quad
\E[X] = \frac{1}{j} \frac{d\Phi_X(\omega)}{d\omega}\bigg|_{\omega=0}.
\]
Plus généralement, pour le moment d'ordre $k$ :
\[
\E[X^k] = \frac{1}{j^k} \frac{d^k \Phi_X(\omega)}{d\omega^k}\bigg|_{\omega=0}.
\]
\end{propbox}

\begin{defnbox}[title=Fonction génératrice des moments]
De la même manière qu'on peut étendre la transformée de Fourier à la transformée de Laplace
en prenant un complexe $s$ au lieu de $j\omega$, on définit la \textbf{fonction génératrice
des moments} :
\[
\hat{\Phi}_X(s) = \E[e^{sX}] = \int_{-\infty}^{\infty} e^{sx} f_X(x)\,dx.
\]
\end{defnbox}

\begin{defnbox}[title=Fonction génératrice de cumulants]
On utilise aussi parfois la \textbf{fonction génératrice de cumulant}, ou \textbf{seconde
fonction caractéristique},
\[
\Psi_X(\omega) = \ln \Phi_X(\omega).
\]
\end{defnbox}

\begin{defnbox}[title=Fonction génératrice de probabilité (cas discret)]
Dans le cas d'une v.a. discrète $X$ prenant des valeurs uniformément espacées $S_X = \{0, 1,
2, 3, \ldots\} = \mathbb{N}$, on définit une \textbf{fonction génératrice de probabilité},
qui est l'espérance de la fonction $z^X$, i.e. la transformée en $z$ de $p_k$ :
\[
G_X(z) = \E[z^X] = \sum_{k=0}^{\infty} z^k P(X=k) = \sum_{k=0}^{\infty} z^k p_k.
\]
\end{defnbox}

\begin{propbox}[title=Propriétés de $G_X(z)$]
On retrouve les probabilités $p_k$ en dérivant $k$ fois $G_X(z)$ par rapport à $z$, et en
évaluant cette dérivée en $z = 0$ :
\[
p_k = P(X=k) = \frac{1}{k!} \frac{d^k G_X(z)}{dz^k}\bigg|_{z=0}.
\]
En dérivant $k$ fois $G_X(z)$ mais en évaluant cette dérivée à présent en $z=1$, on trouve
les moments factoriaux de $X$, à partir desquels on peut retrouver tous les moments de $X$ :
\[
\E[X(X-1)(X-2)\cdots(X-k+1)] = \frac{d^k G_X(z)}{dz^k}\bigg|_{z=1}.
\]
Remarquons finalement que $G_X(1) = \sum_{k=0}^{\infty} p_k = 1$.
\end{propbox}

\subsection{Inégalités}

\begin{rembox}[title=Pourquoi borner plutôt que calculer ?]
Très souvent en pratique, on ne connaît pas complètement la densité de probabilité $f_X(x)$
d'une v.a. $X$. Il est donc très utile de pouvoir borner une probabilité ou l'espérance d'une
fonction de cette v.a.
\end{rembox}

\begin{propbox}[title=Inégalité de Markov]
Si $X \geq 0$ et $a > 0$,
\[
P(X \geq a) \leq \frac{\E[X]}{a}.
\]
\end{propbox}

\begin{propbox}[title=Inégalité de Tchebytcheff]
En posant $D = (X-\mu_X)^2$ et $b = \sqrt{a}$ dans l'inégalité de Markov,
\[
P(|X - \mu_X| \geq b) \leq \frac{\sigma_X^2}{b^2}.
\]
\end{propbox}

\begin{thmbox}[title=Inégalité de Jensen]
Une fonction $g : I \subseteq \mathbb{R} \to \mathbb{R} \cup \{\infty\}$ est \textbf{convexe}
si pour tout $x, x' \in I$ et pour tout $0 \leq t \leq 1$,
\[
g(tx + (1-t)x') \leq t\,g(x) + (1-t)\,g(x').
\]
Une fonction $g$ convexe et continue est toujours minorée par une fonction linéaire passant
par un point de son graphe : pour tout $a \in \mathbb{R}$, il existe $\theta \in \mathbb{R}$
tel que pour tout $x \in \mathbb{R}$, $g(x) \geq g(a) + \theta(x-a)$. En prenant $a = \E[X]$
et $x = X$, puis en prenant l'espérance des deux côtés, il s'ensuit que
\[
\E[g(X)] \geq g(\E[X]).
\]
\end{thmbox}

\end{document}
